\documentclass[11pt]{beamer}
\usetheme{Madrid}
\usefonttheme{serif}

\usepackage[utf8]{inputenc}
\usepackage[spanish]{babel}
\usepackage[T1]{fontenc}

\usepackage{amsmath}
\usepackage{amsfonts}
\usepackage{amssymb}
\usepackage{graphicx, xcolor}
\usepackage{listings}

\usepackage[absolute,overlay]{textpos}
\usepackage[most]{tcolorbox}

\usepackage{tikz}
\usetikzlibrary{shapes.geometric, arrows.meta, positioning, calc}

\tikzstyle{startstop} = [rectangle, rounded corners, minimum width=3cm, minimum height=1cm,text centered, draw=black, fill=red!30]
\tikzstyle{process} = [rectangle, minimum width=3cm, minimum height=1cm, text centered, draw=black, fill=blue!20]
\tikzstyle{decision} = [diamond, aspect=2, minimum width=3.5cm, minimum height=1cm, text centered, draw=black, fill=yellow!40]
\tikzstyle{arrow} = [thick, ->, >=Stealth]


\DeclareMathOperator{\sen}{sen}
\DeclareMathOperator{\tg}{tg}

\setbeamertemplate{caption}[numbered]

%\setbeameroption{show notes on second screen=right}  % También puedes usar: show notes on second screen=right

\author[LLO]{Ing. Leonardo Luis Ortiz}
\title{Simulador HALCON}
% Informe su email de contato o comando a seguir
% Por ejemplo, alcebiades.col@ufes.br
\newcommand{\email}{leonardo.ortiz@ib.edu.ar}
%\setbeamercovered{transparent} 
\setbeamertemplate{navigation symbols}{} 
\logo{\includegraphics[scale=0.08]{images/IB.png}} 
\institute[]{Instituto Balseiro \par Centro Atómico Bariloche} 
\date{} 
%\subject{}

% ---------------------------------------------------------
% Selecione un estilo de referencia
\bibliographystyle{apalike}

%\bibliographystyle{abbrv}
%\setbeamertemplate{bibliography item}{\insertbiblabel}
% ---------------------------------------------------------

% ---------------------------------------------------------
% Incluir los slides 
%\usepackage[spanish,hyperpageref]{backref}

%\renewcommand{\backrefpagesname}{Cita de página(s):~}
%\renewcommand{\backref}{}
%\renewcommand*{\backrefalt}[4]{
	%	\ifcase #1 %
	%		Sin citas en el texto.%
	%	\or
	%		Citado en la página #2.%
	%	\else
	%		Citado #1 veces en las páginas #2.%
	%	\fi}%
% ---------------------------------------------------------
%% Configuración para el código
%\lstset{
%	language=C++,
%	basicstyle=\ttfamily\scriptsize,
%	keywordstyle=\color{blue},
%	commentstyle=\color{gray},
%	stringstyle=\color{orange},
%	showstringspaces=false,
%	breaklines=true,
%	frame=single,
%	numbers=left,
%	numberstyle=\tiny,
%	tabsize=2
%}

% Estilo para VHDL
\lstdefinelanguage{VHDL}{
	morekeywords={
		library,use,all,entity,is,port,in,out,end,architecture,of,begin,
		signal,process,if,then,else,elsif,wait,until,when,others
	},
	sensitive=true,
	morecomment=[l]--,
	morestring=[b]"
}

\lstset{
	language=VHDL,
	basicstyle=\ttfamily\scriptsize,
	keywordstyle=\color{blue},
	commentstyle=\color{gray},
	stringstyle=\color{orange},
	showstringspaces=false,
	breaklines=true,
	frame=single,
	numbers=left,
	numberstyle=\tiny,
	tabsize=2,
	captionpos=b
}


\begin{document}
	
	\begin{frame}
		\titlepage
	\end{frame}
	
	\begin{frame}{Contenido}
		\tableofcontents 
	\end{frame}
	
	\section{Simulador HALCON}
	
	\begin{frame}{HALCON}
		\begin{figure}[htb]
			\centering
			\includegraphics[width=0.9\textwidth]{images/2_HALCON_logo}
		\end{figure}
	\end{frame}
	
	\begin{frame}{HALCON}
		\textbf{HALCON} es un proyecto de código abierto desarrollado por Patricio Reus Merlo que consiste en un conjunto de herramientas desarrolladas en C++20 y Python3 que permiten a los usuarios diseñar y verificar arquitecturas de procesamiento digital de señales (DSP) desde las primeras etapas de diseño (alto nivel) hasta la implementación digital (RTL).
	\end{frame}
	
	\begin{frame}{HALCON}
	\textbf{Entorno de desarrollo}\newline
		
	\begin{itemize}
	\item Lenguaje de programación: C++20
	\item Compilación: Cmake+gcc13 + Python
	\item Control de versiones: GitLab
	\item Entorno de simulación y test: Docker
	\item Automatización y scripting: Python
	\item Documentación de código: Doxygen
	\item Documentación funcional y tutoriales: Markdown+ LaTex
	\end{itemize}
	\end{frame}
	
	\begin{frame}{HALCON - Core}
		\textbf{Estructura de HALCON}
		\begin{figure}[htb]
			\centering
			\includegraphics[width=0.9\textwidth]{images/2_HALCON_structure}
		\end{figure}
	\end{frame}
	
	
	\begin{frame}{HALCON - Core}
	\textbf{Estructura de HALCON}
	\begin{figure}[htb]
		\centering
		\includegraphics[width=1.0\textwidth]{images/2_HALCON_core}
	\end{figure}
	\end{frame}
	
	\begin{frame}{HALCON - Clases de usuario}
	\textcolor{red}{\textbf{Module:}}
	\begin{itemize}
		\item Un módulo es un \underline{contenedor} donde se describen algoritmos y/o modelos
		\item Puede describir estrictamente un bloque de \underline{DSP a nivel de RTL y a nivel de sistemas}
		\item Puede contener modelos de componentes analógicos (ej. canal, ADC)
		\item Puede \underline{instanciar otros módulos} en su interior (submódulos) y configurarlos
		\item Puede \underline{modelar lógica combinacional y secuencial}
	\end{itemize}
	\begin{figure}[htb]
		\centering
		\includegraphics[width=0.8\textwidth]{images/2_HALCON_module}
	\end{figure}
	\end{frame}

	\begin{frame}{HALCON - Clases de usuario}
	\textcolor{red}{\textbf{Module:}}
	\begin{figure}[htb]
		\centering
		\includegraphics[width=0.8\textwidth]{images/2_HALCON_module_2}
	\end{figure}
	\end{frame}
	
	\begin{frame}{HALCON - Clases de usuario}
		\textcolor{red}{\textbf{Register:}}
		\begin{itemize}
			\item Modela un registro de forma explícita y segura (evita registros implícitos)
		\end{itemize}
		\textcolor{red}{\textbf{Port:}}
		\begin{itemize}
			\item La interfase de un módulo con otros módulos se hace a través de puertos
			\item Según su comportamiento frente a un clock, se clasifican en secuenciales o combinacionales
			\item Según el flujo de la información, se clasifican como entradas o salidas
		\end{itemize}
		\begin{figure}[htb]
			\centering
			\includegraphics[width=0.8\textwidth]{images/2_HALCON_registers_ports}
		\end{figure}
	\end{frame}
	
	\begin{frame}{HALCON - Clases de usuario}
		\textcolor{red}{\textbf{Ports y Register:}}
		
		\begin{figure}[htb]
			\centering
			\includegraphics[width=1.0\textwidth]{images/2_HALCON_register_ports_2}
		\end{figure}
	\end{frame}
	
	\begin{frame}{HALCON - Clases de usuario}
	
	\begin{columns}[t]
	\column{0.5\textwidth}	
	\textbf{Puerto combinacioal}
	\begin{itemize}
		\item Rastrean y ejecutan las funciones asociadas a ellos
		\item Aseguran que los valores de los datos estén actualizados
	\end{itemize}	
	\vspace{1.0cm}	
	\begin{figure}[htb]
		\centering
		\includegraphics[width=1.0\textwidth]{images/2_HALCON_puerto_combinacional}
	\end{figure}
		
	\column{0.5\textwidth}	
	\textbf{Puerto secuencial}	
	\begin{itemize}
		\item Lógica controlada por las secuencias de ejecución de reloj
		\item Los valores de los datos dependen de los registros asociados
	\end{itemize}	
	\begin{figure}[htb]
		\centering
		\includegraphics[width=1.0\textwidth]{images/2_HALCON_puerto_secuencial}
	\end{figure}
	
	\end{columns}	
	\end{frame}
	
	\begin{frame}{HALCON - Clases de usuario}
		\textcolor{red}{\textbf{Clock:}}
		
		\begin{itemize}
		\item Contiene un \underline{listado de módulos con sus respectivos registros} que son sensibles a alguno de sus flancos
		\item \underline{Ejecuta la lógica de tiempo “continuo” y de tiempo “discreto”} de los módulos asociados a él
		\item Lleva el \underline{conteo de tiempo} de simulación
		\end{itemize}
		
		\begin{figure}[htb]
			\centering
			\includegraphics[width=0.9\textwidth]{images/2_HALCON_clock}
		\end{figure}
	\end{frame}
	
	\begin{frame}{HALCON - Clases de usuario}
		\textcolor{red}{\textbf{Clock:}}
		\begin{figure}[htb]
			\centering
			\includegraphics[width=0.9\textwidth]{images/2_HALCON_clock_2}
		\end{figure}
	\end{frame}
	
	\begin{frame}{HALCON - Clases de usuario}
		\textcolor{red}{\textbf{Clock:}}
		\begin{figure}[htb]
			\centering
			\includegraphics[width=0.9\textwidth]{images/2_HALCON_clock_3}
		\end{figure}
	\end{frame}
	
	\begin{frame}{HALCON - Clases de usuario}
		\textcolor{red}{\textbf{Clock:}}
		\begin{figure}[htb]
			\centering
			\includegraphics[width=0.9\textwidth]{images/2_HALCON_clock_4}
		\end{figure}
	\end{frame}
	
	\begin{frame}{HALCON - Clases de usuario}
		\textbf{Clocks derivados}

		\begin{figure}[htb]
			\centering
			\includegraphics[width=1.0\textwidth]{images/2_HALCON_clock_5}
		\end{figure}
	\end{frame}
	
	\begin{frame}{HALCON - Clases de usuario}
	\begin{columns}[t]
	\column{0.5\textwidth}	
	\textbf{Clocks derivados}
	\tiny %scriptsize
	\begin{itemize}
	\item Todos los relojes tienen 4 puertos:
	\begin{itemize}
	\tiny
	\item \textcolor{red}{i\_frequency\_hz}: control de frecuencia
	\item \textcolor{red}{i\_phase\_deg}: control de fase inicial
	\item \textcolor{red}{i\_division\_factor\_num}: num. del divisor de frecuencia
	\item \textcolor{red}{i\_division\_factor\_den}: den. del divisor de frecuencia
	\end{itemize}
	
	\item El puerto \textcolor{red}{i\_frequency\_hz} de los clocks derivados se conecta automáticamente al puerto i\_frequency\_hz del clock del que se derivan.
	\item \textcolor{red}{IMPORTANTE}: luego de la derivación de un clock, no redefinir el puerto \textcolor{red}{i\_frequency\_hz} del clock derivado.
	\item 	Cada vez que se ejecuta un clock, éste se actualiza leyendo el valor de frecuencia de su puerto. A este puerto se puede conectar una señal variante en el tiempo.
	\item Todas las conexiones entre clocks y las derivaciones se deben hacer dentro de un método Connect() del módulo/simulador que contenga a los clocks.
	
	\end{itemize}	

	\column{0.5\textwidth}	
	\begin{figure}[htb]
		\centering
		\includegraphics[width=1.0\textwidth]{images/2_HALCON_clock_6}
	\end{figure}
	\end{columns}	
	\end{frame}
	
		\begin{frame}{HALCON - Clases de usuario}
		\textcolor{red}{\textbf{Clock:}}
		\begin{figure}[htb]
			\centering
			\includegraphics[width=0.9\textwidth]{images/2_HALCON_clock_2}
		\end{figure}
	\end{frame}
	
	\begin{frame}{HALCON - Clases de usuario}
		\textcolor{red}{\textbf{Simulator:}}
		\begin{itemize}
			\item Al derivar de módulo hereda todas sus propiedades. Es decir, es un contenedor de otros bloques (ej. módulos y clocks)
			\item Agrega bloques básicos que hacen evolucionar la simulación: Scheduler, Command Handler, Logger, Setter, etc
			\item Es el bloque de mayor jerarquía y por ende es el encargado de nombrar, configurar, conectar e inicializar todos los módulos y submódulos que contenga
		\end{itemize}
		\begin{figure}[htb]
			\centering
			\includegraphics[width=0.9\textwidth]{images/2_HALCON_simulator}
		\end{figure}
	\end{frame}
	
	
		\begin{frame}{HALCON - Clases de usuario}
		\textcolor{red}{\textbf{Simulator:}}
		\begin{itemize}
			\item Al derivar de módulo hereda todas sus propiedades. Es decir, es un contenedor de otros bloques (ej. módulos y clocks)
			\item Agrega bloques básicos que hacen evolucionar la simulación: Scheduler, Command Handler, Logger, Setter, etc
			\item Es el bloque de mayor jerarquía y por ende es el encargado de nombrar, configurar, conectar e inicializar todos los módulos y submódulos que contenga
		\end{itemize}
		\begin{figure}[htb]
			\centering
			\includegraphics[width=0.9\textwidth]{images/2_HALCON_simulator}
		\end{figure}
	\end{frame}
	
	\lstdefinestyle{myCustomMatlabStyle}{
		language=Matlab,
		numbers=left,
		stepnumber=1,
		numbersep=10pt,
		tabsize=4,
		showspaces=false,
		showstringspaces=false
	}
	
	\begin{frame}{\textbf{HALCON} - Startup}
	\textcolor{red}{\textbf{Simulator:}}
	\begin{figure}[htb]
	\begin{flushright} 
		\includegraphics[width=0.6\textwidth]{images/2_HALCON_simulator_1_code}
	\end{flushright}
	\end{figure}
	\begin{figure}[htb]
		\centering
		\includegraphics[width=0.9\textwidth]{images/2_HALCON_simulator_1}
	\end{figure}
	\end{frame}


	\begin{frame}{\textbf{HALCON} - Startup}	
	\textcolor{red}{\textbf{Simulator:}}
	\begin{figure}[htb]
	\begin{flushright} 
	\includegraphics[width=0.45\textwidth]{images/2_HALCON_simulator_2_code}
	\end{flushright}
	\end{figure}
	
	\begin{figure}[htb]
		\centering
		\includegraphics[width=0.8\textwidth]{images/2_HALCON_simulator_2}
	\end{figure}
	\end{frame}


	\begin{frame}{\textbf{HALCON} - Startup}
	\textcolor{red}{\textbf{Simulator:}}
	\begin{figure}[htb]
		\begin{flushright} 
			\includegraphics[width=0.45\textwidth]{images/2_HALCON_simulator_3_code}
		\end{flushright}
	\end{figure}
	
	\begin{figure}[htb]
		\centering
		\includegraphics[width=0.8\textwidth]{images/2_HALCON_simulator_3}
	\end{figure}
	\end{frame}

	\begin{frame}{\textbf{HALCON} - Startup}
	\textcolor{red}{\textbf{Simulator:}}
	\begin{figure}[htb]
		\begin{flushright} 
			\includegraphics[width=0.6\textwidth]{images/2_HALCON_simulator_4_code}
		\end{flushright}
	\end{figure}
	
	\begin{figure}[htb]
		\centering
		\includegraphics[width=0.8\textwidth]{images/2_HALCON_simulator_4}
	\end{figure}
	\end{frame}

	\begin{frame}{\textbf{HALCON} - Startup}
	\textcolor{red}{\textbf{Simulator:}}
	\begin{figure}[htb]
		\begin{flushright} 
			\includegraphics[width=0.55\textwidth]{images/2_HALCON_simulator_5_code}
		\end{flushright}
	\end{figure}
	
	\begin{figure}[htb]
		\centering
		\includegraphics[width=0.8\textwidth]{images/2_HALCON_simulator_5}
	\end{figure}
	\end{frame}
	
	
	\begin{frame}{\textbf{HALCON} - Startup}
	\textcolor{red}{\textbf{Simulator:}}
	\begin{figure}[htb]
		\begin{flushright} 
			\includegraphics[width=0.55\textwidth]{images/2_HALCON_simulator_6_code}
		\end{flushright}
	\end{figure}
	
	\begin{figure}[htb]
		\centering
		\includegraphics[width=0.8\textwidth]{images/2_HALCON_simulator_6}
	\end{figure}
	\end{frame}
	
	
	\begin{frame}{\textbf{HALCON} - Startup}
	\textcolor{red}{\textbf{Simulator:}}
	\begin{figure}[htb]
		\begin{flushright} 
			\includegraphics[width=0.55\textwidth]{images/2_HALCON_simulator_7_code}
		\end{flushright}
	\end{figure}
	
	\begin{figure}[htb]
		\centering
		\includegraphics[width=0.8\textwidth]{images/2_HALCON_simulator_7}
	\end{figure}
	\end{frame}
	
	
	\begin{frame}{\textbf{HALCON} - Startup}
	\textcolor{red}{\textbf{Simulator:}}
	\begin{figure}[htb]
		\begin{flushright} 
			\includegraphics[width=0.55\textwidth]{images/2_HALCON_simulator_8_code}
		\end{flushright}
	\end{figure}
	
	\begin{figure}[htb]
		\centering
		\includegraphics[width=0.8\textwidth]{images/2_HALCON_simulator_8}
	\end{figure}
	\end{frame}

	\begin{frame}{\textbf{HALCON} - Startup}
	\textcolor{red}{\textbf{Simulator:}}
	\begin{figure}[htb]
		\begin{flushright} 
			\includegraphics[width=0.55\textwidth]{images/2_HALCON_simulator_9_code}
		\end{flushright}
	\end{figure}
	
	\begin{figure}[htb]
		\centering
		\includegraphics[width=0.8\textwidth]{images/2_HALCON_simulator_9}
	\end{figure}
	\end{frame}
	
	\begin{frame}{\textbf{HALCON} - Startup}
	\textcolor{red}{\textbf{Simulator:}}

	\begin{figure}[htb]
		\centering
		\includegraphics[width=1.0\textwidth]{images/2_HALCON_simulator_10}
	\end{figure}
	\end{frame}
	
	\section{Tools}

	\begin{frame}{\textbf{HALCON} - LOGs}
	\begin{columns}[t]
		\column{0.5\textwidth}	
		\scriptsize
		\begin{itemize}
		\item \textcolor{red}{signals}: señal de interés
		\item \textcolor{red}{clock}: reloj de muestreo
		\item \textcolor{red}{edge}: flancos de muestreo. Puede ser POSITIVE, NEGATIVE o BOTH
		\item \textcolor{red}{begin, step, end}: definen la ventana de muestreo y el subsample
		\item \textcolor{red}{file\_type}: define el tipo de archive de salida. Puede ser TEXT o BINARY
		\item \textcolor{red}{file\_name\_type}: controla el nombre del archive. Puede ser SHORT o LONG
		\item \textcolor{red}{format}: solo afecta a los LOGs de datos de tipo \underline{ac\_fixed}, \underline{array ac\_fixed}, \underline{Port ac\_fixed} y \underline{Port array ac\_fixed}. Solo válido con archivos de texto. Las macros soportadas son DOUBLE, UINT y HEX
		\end{itemize}	
		
		\column{0.5\textwidth}	
		\begin{figure}[htb]
			\centering
			\includegraphics[width=1.0\textwidth]{images/2_HALCON_LOGS}
		\end{figure}
	\end{columns}	
	\end{frame}
	
	\begin{frame}{\textbf{HALCON} - LOGs}
	\begin{figure}[htb]
	\centering
	\includegraphics[width=1.0\textwidth]{images/2_HALCON_LOGS_2}
	\end{figure}
	\end{frame}
	
	\begin{frame}{\textbf{HALCON}}
	\textbf{Barrido de parámetros: clases}	
	\begin{figure}[htb]
		\begin{flushleft}
		\includegraphics[width=0.55\textwidth]{images/2_HALCON_sweep}
		\end{flushleft}
	\end{figure}
	\end{frame}
	
	\begin{frame}{\textbf{HALCON}}
	\textbf{Barrido de parámetros: clases}	
	\begin{figure}[htb]
		\centering
		\includegraphics[width=1.0\textwidth]{images/2_HALCON_sweep_2}
	\end{figure}
	\end{frame}
	
	\begin{frame}{\textbf{HALCON}}
	\textbf{Barrido de parámetros: clases}	
	\begin{figure}[htb]
		\centering
		\includegraphics[width=1.0\textwidth]{images/2_HALCON_sweep_3}
	\end{figure}
	\end{frame}
	
	\begin{frame}{\textbf{HALCON}}
	\textbf{Barrido de parámetros: clases}	
	\begin{figure}[htb]
		\centering
		\includegraphics[width=1.0\textwidth]{images/2_HALCON_sweep_4}
	\end{figure}
	\end{frame}
	
	\begin{frame}{\textbf{HALCON}}
	\textbf{Barrido de parámetros: directorios}	
	\begin{figure}[htb]
		\centering
		\includegraphics[width=1.0\textwidth]{images/2_HALCON_sweep_5}
	\end{figure}
	\end{frame}
	
	\begin{frame}{\textbf{HALCON}}
	\textbf{Barrido de parámetros: clases}	
	\begin{figure}[htb]
		\centering
		\includegraphics[width=1.0\textwidth]{images/2_HALCON_sweep_6}
	\end{figure}
	\end{frame}
	
	\begin{frame}{\textbf{HALCON}}
	\textbf{Barrido de parámetros: clases}	
	\begin{figure}[htb]
		\centering
		\includegraphics[width=1.0\textwidth]{images/2_HALCON_sweep_7}
	\end{figure}
	\end{frame}

	\section{Ejemplo}
	
	\begin{frame}
		\centering
		\Huge \textbf{Ejemplo}
	\end{frame}
	
	\begin{frame}{\textbf{HALCON} - Registro de desplazamiento}
	\begin{columns}[t]
		\column{0.5\textwidth}	
		\begin{figure}[htb]
			\centering
			\includegraphics[width=1.0\textwidth]{images/2_HALCON_sr_1}
		\end{figure}
		
		\column{0.5\textwidth}	
		\begin{figure}[htb]
			\centering
			\includegraphics[width=1.0\textwidth]{images/2_HALCON_sr_1_code}
		\end{figure}
	\end{columns}	
	\end{frame}
	
	\begin{frame}{\textbf{HALCON} - Registro de desplazamiento}
		\begin{columns}[t]
			\column{0.5\textwidth}	
			\begin{figure}[htb]
				\centering
				\includegraphics[width=1.0\textwidth]{images/2_HALCON_sr_1}
			\end{figure}
			
			\column{0.5\textwidth}	
			\begin{figure}[htb]
				\centering
				\includegraphics[width=1.0\textwidth]{images/2_HALCON_sr_2_code}
			\end{figure}
		\end{columns}	
	\end{frame}
	
	\begin{frame}{\textbf{HALCON} - Registro de desplazamiento}
		\begin{columns}[t]
			\column{0.5\textwidth}	
			\begin{figure}[htb]
				\centering
				\includegraphics[width=1.0\textwidth]{images/2_HALCON_sr_1}
			\end{figure}
			
			\column{0.5\textwidth}	
			\begin{figure}[htb]
				\centering
				\includegraphics[width=1.0\textwidth]{images/2_HALCON_sr_3_code}
			\end{figure}
		\end{columns}	
	\end{frame}
	
	% Testbench
	
	\begin{frame}{\textbf{HALCON} - Registro de desplazamiento}
		\begin{columns}[t]
			\column{0.5\textwidth}	
			\begin{figure}[htb]
				\centering
				\includegraphics[width=1.0\textwidth]{images/2_HALCON_sr_4}
			\end{figure}
			
			\column{0.5\textwidth}	
			\begin{figure}[htb]
				\centering
				\includegraphics[width=1.0\textwidth]{images/2_HALCON_sr_4_code}
			\end{figure}
		\end{columns}	
	\end{frame}
	
	\begin{frame}{\textbf{HALCON} - Registro de desplazamiento}
		\begin{columns}[t]
			\column{0.5\textwidth}	
			\begin{figure}[htb]
				\centering
				\includegraphics[width=1.0\textwidth]{images/2_HALCON_sr_5_code_1}
			\end{figure}
			
			\column{0.5\textwidth}	
			\begin{figure}[htb]
				\centering
				\includegraphics[width=1.0\textwidth]{images/2_HALCON_sr_5_code_2}
			\end{figure}
		\end{columns}	
	\end{frame}
	
	\begin{frame}{\textbf{HALCON} - Registro de desplazamiento}
		\begin{columns}[t]
			\column{0.5\textwidth}	
			\begin{figure}[htb]
				\centering
				\includegraphics[width=1.0\textwidth]{images/2_HALCON_sr_6_code_1}
			\end{figure}
			
			\column{0.5\textwidth}	
			\begin{figure}[htb]
				\centering
				\includegraphics[width=1.0\textwidth]{images/2_HALCON_sr_6_code_2}
			\end{figure}
		\end{columns}	
	\end{frame}
	
	\begin{frame}{\textbf{HALCON} - Registro de desplazamiento}
		\begin{columns}[t]
			\column{0.45\textwidth}	
			\begin{figure}[htb]
				\centering
				\includegraphics[width=1.0\textwidth]{images/2_HALCON_sr_7_code_1}
			\end{figure}
			
			\column{0.55\textwidth}	
			\begin{figure}[htb]
				\centering
				\includegraphics[width=1.0\textwidth]{images/2_HALCON_sr_7_code_2}
			\end{figure}
		\end{columns}	
	\end{frame}
	
	
	
	
	
	
	
	
	
	\begin{frame}
	\centering
	\Huge \textbf{Q\&A}
	\end{frame}
	
%	\begin{frame}{Lenguajes de Descripción de Hardware - HDL}
%		
%	\end{frame}
%	
%	
%	\begin{frame}[fragile]{Lenguajes de Descripción de Hardware - VHDL}
%
%	\end{frame}
%	
%	\begin{frame}[fragile]{Lenguajes de Descripción de Hardware - VHDL}
%
%	\end{frame}
%	
%	\begin{frame}[fragile]{Lenguajes de Descripción de Hardware - VHDL Testbench)}
%
%	\end{frame}
%
%	\begin{frame}[fragile]{Lenguajes de Descripción de Hardware - VHDL Testbench}
%
%	\end{frame}
%
%	\begin{frame}[fragile]{Lenguajes de Descripción de Hardware - VHDL Testbench}
%
%	\end{frame}
%
%	\begin{frame}[fragile]{Lenguajes de Descripción de Hardware - VHDL Testbench}
%
%	\end{frame}
%
%
%\begin{frame}{Lenguajes de Descripción de Hardware - Resumen}
%	
%\end{frame}
%
%\begin{frame}{Lenguajes de Descripción de Hardware - Resumen}
%
%\end{frame}
%
%\begin{frame}{Lenguajes de Descripción de Hardware - Resumen}	
%
%\end{frame}
%
%
%\begin{frame}{Lenguajes de Software - Resumen}
%
%\end{frame}
%
%\begin{frame}{Lenguajes de Software - Resumen}
%
%\end{frame}
%
%\begin{frame}{Puntos claves: Concurrency}
%
%\end{frame}
%
%\begin{frame}{Puntos claves: Concurrency}
%
%\end{frame}
%
%\begin{frame}{Puntos claves: Concurrency}
%
%\end{frame}
%
%\begin{frame}{Puntos claves: Timing}
%
%\end{frame}
%
%\begin{frame}{Puntos claves: Timing}
%
%\end{frame}
%
%\begin{frame}{Puntos claves: Applications}
%
%\end{frame}
%
%\begin{frame}{Puntos claves: Applications}
%
%\end{frame}
%
%\section{Programación controlada por eventos}
%	\begin{frame}{Programación controlada por eventos}
%	\end{frame}
%	
%
%\begin{frame}{Programación controlada por eventos}
%
%\end{frame}	
%	
%
%\begin{frame}{Programación controlada por eventos}
%
%\end{frame}	
%
%
%\begin{frame}[fragile]{Programación controlada por eventos}
%
%\end{frame}


%	
%\section{Utilizando C++ y Python}
%	\begin{frame}{Utilizando C++ y Python}
%		Breve descripción de Lenguajes de Descripción de Hardware. Comparativa
%	\end{frame}
%	
%\section{Simulador HALCON - Getting started}
%	\begin{frame}{Simulador HALCON - Getting started}
%		Breve descripción de Lenguajes de Descripción de Hardware. Comparativa
%	\end{frame}
%	
%\section{Ejemplo 1}
%	\begin{frame}{Ejemplo 1}
%		Breve descripción de Lenguajes de Descripción de Hardware. Comparativa
%	\end{frame}
%	
%\section{Ejemplo 2}
%	\begin{frame}{Ejemplo 2}
%		Breve descripción de Lenguajes de Descripción de Hardware. Comparativa
%	\end{frame}
%
%\section{Ejemplo 3}
%	\begin{frame}{Ejemplo 3}
%		Breve descripción de Lenguajes de Descripción de Hardware. Comparativa
%	\end{frame}
%
%\section{Ejemplo 4}
%	\begin{frame}{Ejemplo 4}
%		Breve descripción de Lenguajes de Descripción de Hardware. Comparativa
%	\end{frame}
%
%\section{Trabajo Final. Propuestas e Inicio implementación.}
%\section{Trabajo Final. Implementación}
%	
%	
%%
%%	
%%	\begin{frame}{Definiciones, Teoremas, Corolarios e Ejemplos}
%%		\begin{block}{Definición 2.1.1}
%%			Aquí tenemos una nueva definición
%%		\end{block}
%%		
%%		\begin{block}{Teorema 2.1.2}
%%			Aquí tenemos un nuevo teorema
%%		\end{block}
%%		
%%		\begin{block}{Ejemplo 2.1.3}
%%			Aquí tenemos un nuevo ejemplo.
%%		\end{block}
%%		
%%		\begin{block}{Corolario 2.1.4}
%%			Aqui tenemos un nuevo corolario.
%%		\end{block}
%%	\end{frame}
%%	
%%	\begin{frame}{Citas de las referencias}
%%		Libro de Análisis Real~\cite{lima2004analise}
%%	\end{frame}
%%	
%%	\begin{frame}{Listas}
%%		Ejemplo de lista con símbolos:
%%		
%%		\begin{itemize}
%%			\item Primer elemento
%%			\item Segundo elemento
%%		\end{itemize}
%%		
%%		\bigskip
%%		
%%		Ejemplo de lista con items enumerados:
%%		
%%		\begin{enumerate}
%%			\item Primer elemento
%%			\item Segundo elemento
%%		\end{enumerate}
%%	\end{frame}
%%	
%%	\begin{frame}{Tablas}
%%		A Tabla muestra un modelo de tabla.
%%		
%%		\begin{table}[htb]
%%			\caption{Ejemplo de tabla}
%%			\label{tab:modelo_tabla}
%%			\centering
%%			\begin{tabular}{c|c|c|c}
%%				\hline
%%				\textbf{Persona} & \textbf{Edad} & \textbf{Peso} & \textbf{Altura} \\ \hline
%%				Marcos & 26    & 68   & 178    \\ \hline
%%				Ivone  & 22    & 57   & 162    \\ \hline
%%				Sueli  & 40    & 65   & 153    \\ \hline
%%			\end{tabular}
%%			
%%			\medskip
%%			
%%			Fuente: Producción del propio autor.
%%		\end{table}
%%	\end{frame}
%%	
%%	\begin{frame}{Imagenes}
%%		La Figura muestra el logomarca promocional de Universidade Federal do Espírito Santo (UFES).
%%		
%%		\begin{figure}[htb]
%%			\caption{Logomarca Promocional UFES}
%%			\label{fig:logo_UFES}
%%			\centering
%%			\includegraphics[width=0.5\textwidth]{imagens/logomarca_UFES}
%%			
%%			\medskip
%%			
%%			Fuente: Producción del propio autor.
%%		\end{figure}
%%	\end{frame}
%%	
%%	\begin{frame}{Conclusión}
%%		Texto referente a la conclusión.
%%	\end{frame}
%%	
%%	\begin{frame}{Trabajos a futuro}
%%		\begin{itemize}
%%			\item Estudiar $\ldots$
%%			
%%			\medskip
%%			
%%			\item Explorar $\ldots$
%%			
%%			\medskip
%%			
%%			\item Analizar $\ldots$
%%		\end{itemize}
%%	\end{frame}
%%	
%%	\section{Referencias}
%%	\begin{frame}[allowframebreaks]{Referencias}
%%		\bibliography{referencias}
%%	\end{frame}
%%	
%%	\begin{frame}
%%		
%%		\begin{center}
%%			Gracias!
%%			
%%			\email
%%		\end{center}
%%		
%%		\begin{figure}[htb]
%%			\centering
%%			\includegraphics[width=0.5\textwidth]{imagens/IB}
%%		\end{figure}
%%		
%%	\end{frame}
	
\end{document}